
\documentclass{article}
\usepackage{colortbl}
\usepackage{makecell}
\usepackage{multirow}
\usepackage{supertabular}

\begin{document}

\newcounter{utterance}

\centering \large Interaction Transcript for game `wordle', experiment `high\_frequency\_words\_no\_clue\_no\_critic', episode 5 with claude{-}opus{-}4{-}8{-}azure.
\vspace{24pt}

{ \footnotesize  \setcounter{utterance}{1}
\setlength{\tabcolsep}{0pt}
\begin{supertabular}{c@{$\;$}|p{.15\linewidth}@{}p{.15\linewidth}p{.15\linewidth}p{.15\linewidth}p{.15\linewidth}p{.15\linewidth}}
    \# & \multicolumn{2}{c}{Player} && \multicolumn{2}{c}{Game Master} \\
    \hline

    \theutterance \stepcounter{utterance}  
    & & & \multicolumn{4}{p{0.6\linewidth}}{
        \cellcolor[rgb]{0.9,0.9,0.9}{
            \makecell[{{p{\linewidth}}}]{
                \texttt{\tiny{[P1$\langle$GM]}}
                \texttt{Jesteś czarodziejem języka, który lubi zgadywać słowa, korzystając z podanych zasad.} \\
\\ 
\texttt{Witamy w Wordle! Masz sześć prób, aby odgadnąć słowo docelowe, poprawne polskie słowo składające się z pięciu małych liter. Użyj tagów "wyjaśnienie:" oraz "zgadnij:", aby podać zwięzłe wyjaśnienie dla każdej próby.} \\
\\ 
\texttt{Na przykład, jeśli Twoją próbą jest "dzień", Twoja odpowiedź powinna brzmieć} \\
\texttt{wyjaśnienie: to jest popularne pięcioliterowe polskie słowo i zaczynam zgadywanie od tego słowa.} \\
\texttt{zgadnij: dzień} \\
\\ 
\texttt{Po każdej próbie Twoja odpowiedź zostanie zweryfikowana i otrzymasz informację zwrotną wskazującą, które litery są poprawne (zielony), które litery są poprawne, ale na niewłaściwej pozycji (żółty), oraz które litery są niepoprawne (czerwony). Ta informacja zwrotna może być przydatna przy ustalaniu, które litery uwzględnić lub wykluczyć w następnej próbie.} \\
\\ 
\texttt{Na przykład informacja zwrotna dla "dzień" może wyglądać tak:} \\
\texttt{informacja zwrotna: d<żółty> z<żółty> i<zielony> e<żółty> ń<czerwony>} \\
\\ 
\texttt{Wyjaśnienie powinno zawierać szczegóły dotyczące tego, jak wykorzystano informacja zwrotna, aby dojść do nowej próby.} \\
\\ 
\texttt{Nie generuj żadnego innego tekstu.} \\
\\ 
\texttt{Zacznijmy od Twojej pierwszej próby.} \\
            }
        }
    }
     \\ \\

    \theutterance \stepcounter{utterance}  
    & \multicolumn{4}{p{0.6\linewidth}}{
        \cellcolor[rgb]{0.561,0.737,0.561}{
            \makecell[{{p{\linewidth}}}]{
                \texttt{\tiny{[P1$\rangle$GM]}}
                \texttt{wyjaśnienie: zaczynam od popularnego pięcioliterowego polskiego słowa z częstymi literami.} \\
\texttt{zgadnij: kotem} \\
            }
        }
    }
    & & \\ \\

    \theutterance \stepcounter{utterance}  
    & & & \multicolumn{2}{p{0.3\linewidth}}{
        \cellcolor[rgb]{0.9,0.9,0.9}{
            \makecell[{{p{\linewidth}}}]{
                \texttt{\tiny{[GM$|$GM]}}
                \texttt{attempts = 1} \\
\texttt{target\_word = karta} \\
\texttt{guess = kotem} \\
\texttt{guess\_feedback = k<green> o<red> t<yellow> e<red> m<red>} \\
            }
        }
    }
    & & \\ \\

    \theutterance \stepcounter{utterance}  
    & & & \multicolumn{4}{p{0.6\linewidth}}{
        \cellcolor[rgb]{0.9,0.9,0.9}{
            \makecell[{{p{\linewidth}}}]{
                \texttt{\tiny{[P1$\langle$GM]}}
                \texttt{informacja zwrotna: k<zielony> o<czerwony> t<żółty> e<czerwony> m<czerwony>} \\
\\ 
\texttt{Podaj odpowiedź tylko w tym formacie.} \\
\texttt{wyjaśnienie: szczegóły} \\
\texttt{zgadnij: słowo} \\
            }
        }
    }
     \\ \\

    \theutterance \stepcounter{utterance}  
    & \multicolumn{4}{p{0.6\linewidth}}{
        \cellcolor[rgb]{0.561,0.737,0.561}{
            \makecell[{{p{\linewidth}}}]{
                \texttt{\tiny{[P1$\rangle$GM]}}
                \texttt{wyjaśnienie: K jest zielone na pozycji 1, T jest w słowie ale nie na pozycji 3. Wykluczam O, E, M. Próbuję słowa z T na innej pozycji.} \\
\texttt{zgadnij: karty} \\
            }
        }
    }
    & & \\ \\

    \theutterance \stepcounter{utterance}  
    & & & \multicolumn{2}{p{0.3\linewidth}}{
        \cellcolor[rgb]{0.9,0.9,0.9}{
            \makecell[{{p{\linewidth}}}]{
                \texttt{\tiny{[GM$|$GM]}}
                \texttt{attempts = 2} \\
\texttt{target\_word = karta} \\
\texttt{guess = karty} \\
\texttt{guess\_feedback = k<green> a<green> r<green> t<green> y<red>} \\
            }
        }
    }
    & & \\ \\

    \theutterance \stepcounter{utterance}  
    & & & \multicolumn{4}{p{0.6\linewidth}}{
        \cellcolor[rgb]{0.9,0.9,0.9}{
            \makecell[{{p{\linewidth}}}]{
                \texttt{\tiny{[P1$\langle$GM]}}
                \texttt{informacja zwrotna: k<zielony> a<zielony> r<zielony> t<zielony> y<czerwony>} \\
\\ 
\texttt{Podaj odpowiedź tylko w tym formacie.} \\
\texttt{wyjaśnienie: szczegóły} \\
\texttt{zgadnij: słowo} \\
            }
        }
    }
     \\ \\

    \theutterance \stepcounter{utterance}  
    & \multicolumn{4}{p{0.6\linewidth}}{
        \cellcolor[rgb]{0.561,0.737,0.561}{
            \makecell[{{p{\linewidth}}}]{
                \texttt{\tiny{[P1$\rangle$GM]}}
                \texttt{wyjaśnienie: K, A, R, T są zielone na pozycjach 1{-}4. Y jest niepoprawne, więc zamieniam ostatnią literę. "Karta" pasuje idealnie.} \\
\texttt{zgadnij: karta} \\
            }
        }
    }
    & & \\ \\

    \theutterance \stepcounter{utterance}  
    & & & \multicolumn{2}{p{0.3\linewidth}}{
        \cellcolor[rgb]{0.9,0.9,0.9}{
            \makecell[{{p{\linewidth}}}]{
                \texttt{\tiny{[GM$|$GM]}}
                \texttt{attempts = 3} \\
\texttt{target\_word = karta} \\
\texttt{guess = karta} \\
\texttt{guess\_feedback = k<green> a<green> r<green> t<green> a<green>} \\
            }
        }
    }
    & & \\ \\

    \theutterance \stepcounter{utterance}  
    & & & \multicolumn{2}{p{0.3\linewidth}}{
        \cellcolor[rgb]{0.9,0.9,0.9}{
            \makecell[{{p{\linewidth}}}]{
                \texttt{\tiny{[GM$|$GM]}}
                \texttt{game\_result = WIN} \\
            }
        }
    }
    & & \\ \\

\end{supertabular}
}

\end{document}
